\item Una expresión general para los niveles de energía de los átomos e iones de un solo electrón es:
$$ E_n = -\frac{\mu k_e^2 q_1^2 q_2^2}{2\hbar^2 n^2} $$
Aquí $\mu$ es la masa reducida, dada por $\mu = m_1 m_2 / (m_1 + m_2)$, donde $m_1$ es la masa del electrón y $m_2$ es la masa del núcleo; $k_e$ es la constante de Coulomb, y $q_1$ y $q_2$ son las cargas del electrón y del núcleo respectivamente. La longitud de onda para la transición $n = 3$ a $n = 2$ en el átomo de hidrógeno es $656.3\text{ nm}$. ¿Cuáles son las longitudes de onda para esta misma transición en:
\begin{enumerate}
    \item positronio (constituido por un electrón y un positrón)?
    \item helio con una sola ionización?
\end{enumerate}